\section{Тест производительности}

Для оценки эффективности Суффиксного дерева (алгоритм Укконена, $O(N)$) было проведено его сравнение с \textbf{наивным алгоритмом} (Brute Force, $O(N^2)$) и узкоспециализированным \textbf{алгоритмом Бута} ($O(N)$). Тестирование проводилось на вырожденных строках, где время работы наивного алгоритма стремится к квадратичному.

Результаты замеров времени и потребления оперативной памяти приведены в таблице \ref{tab:benchmark}.

\begin{table}[h!]
\centering
\caption{Сравнение производительности алгоритмов}
\label{tab:benchmark}
\begin{small}
\begin{tabular}{|l|c|c|c|c|c|c|c|}
\hline
\multirow{2}{*}{\textbf{N символов}} & \multicolumn{2}{c|}{\textbf{Наивный ($O(N^2)$)}} & \multicolumn{2}{c|}{\textbf{Алгоритм Бута}} & \multicolumn{2}{c|}{\textbf{Суфф. дерево}} & \multirow{2}{*}{\textbf{Ускорение*}} \\ \cline{2-7}
 & Время & Память & Время & Память & Время & Память & \\ \hline
500 000 & 41.02 с & 5.3 MB & 0.00 с & 5.4 MB & 0.12 с & 220.7 MB & $\approx 341 \times$ \\ \hline
10 000 000 & ---** & --- & 0.01 с & 42.1 MB & 9.37 с & 4.2 GB & $> 1700 \times$ \\ \hline
\end{tabular}
\end{small}

\smallskip
\begin{flushleft}
\footnotesize
* Ускорение рассчитано как отношение времени наивного алгоритма к времени построения суффиксного дерева. \\
** Тестирование наивного алгоритма на 10 млн символов было прервано из-за превышения лимита времени (ожидаемое время > 4 часов).
\end{flushleft}
\end{table}

\subsection*{Анализ результатов}

1. \textbf{Производительность.} На тесте из $500,000$ символов суффиксное дерево показало ускорение в 341 раз относительно наивного подхода. При увеличении входных данных до $10$ миллионов символов, время работы дерева выросло линейно, в то время как наивный алгоритм не смог завершить работу в разумные сроки. Алгоритм Бута остается самым быстрым за счет отсутствия сложных структур данных.

2. \textbf{Память.} Суффиксное дерево продемонстрировало крайне высокое потребление ресурсов. Для строки в 10 млн символов потребовалось более 4 GB оперативной памяти. Это подтверждает, что дерево Укконена требует больше памяти, чем сама исходная строка, из-за хранения узлов, суффиксных ссылок и словарей \texttt{std::map} для каждого перехода.

3. \textbf{Вывод.} Суффиксное дерево оправдано использовать в задачах, где требуется многократный поиск или сложный анализ текста. Для разовой задачи поиска минимального циклического сдвига эффективнее использовать алгоритм Бута.

\pagebreak